\documentclass[11pt]{article}
\usepackage{lmodern}
\usepackage[T1]{fontenc}
\usepackage[letterpaper,margin=0.8in]{geometry}
\usepackage{xcolor}
\usepackage{booktabs}
\usepackage{tabularx}
\usepackage{longtable}
\usepackage{enumitem}
\usepackage{float}
\usepackage{array}
\usepackage{needspace}
\usepackage[american]{circuitikz}
\usetikzlibrary{positioning,fit,arrows.meta}
\usepackage[hidelinks]{hyperref}

\setlist{itemsep=1pt,topsep=3pt}
\newcommand{\tbd}[1]{\textcolor{red}{[TBD: #1]}}
\newcommand{\opt}[1]{\texttt{-{}-#1}}
\newcolumntype{L}{>{\raggedright\arraybackslash}X}

\title{PACMAN Quality Assurance and Quality Control\\[2pt]
  \large Hardware checkout: tools, coverage, and shifter's guide}
\author{Draft maintained by the PACMAN agents (augie); reference for the TRM}
\date{2026-10-17}

\begin{document}
\maketitle
\sloppy

\section{Purpose and scope}

\begin{itemize}
\item Define a repeatable checkout of every new PACMAN card before it goes into a production system or onto an
  ASIC test stand.
\item Two utilities, both run on the PACMAN's own Linux system:
  \begin{itemize}
  \item \texttt{pacman\_lnxcheck}: the standard embedded-Linux checkout of the Trenz module (CPU, memory,
    on-chip supplies, Ethernet, flash and SD storage).
  \item \texttt{pacman\_hwcheck}: the PACMAN-specific checkout (power, tile I/O, timing inputs, front panel),
    partly guided by an operator.
  \end{itemize}
\item Scope: the features needed for normal operation.  Features that normal operation does not use are listed
  (Section~\ref{sec:status}) but not tested.
\item Method: wherever possible a feature is tested by loopback, so that the feature and its feedback path are
  verified together, and a failure points at hardware, firmware or software.
\end{itemize}

\section{Production setting}

\begin{itemize}
\item PACMAN 1v5 with a Trenz TE0720 module (Zynq-7000).  The card boots from its microSD card (boot image and
  Linux root file system) and is reached over Ethernet only (\texttt{ssh}).  Production cards have no USB/JTAG
  header installed.
\item \textbf{ASIC protection jumper (J5).}  Installed whenever the card is in the production system or connected
  to ASICs.  The checkout runs with J5 \emph{removed}; with J5 installed, \texttt{pacman\_hwcheck} refuses
  everything except its read-only report.
\item The PACMAN data and command servers are stopped during the checkout (they share the data path with the
  loopback test) and restarted afterwards.
\end{itemize}


\section{Test setup}

\begin{figure}[H]
\centering
\resizebox{\linewidth}{!}{%
\begin{circuitikz}[
  box/.style={draw,thick,rounded corners=2pt,align=center,font=\small},
  jack/.style={draw,thick,circle,minimum size=5mm,inner sep=0pt,fill=white},
  port/.style={draw,thick,minimum width=7mm,minimum height=5mm,inner sep=0pt,fill=white},
  sig/.style={font=\scriptsize,anchor=west},
  ublnk/.style={-{Latex[length=2mm]},thick,rounded corners=4pt},
  tbc/.style={font=\scriptsize\itshape,text=gray}]
  \ctikzset{bipoles/length=0.9cm}
  % ---- PACMAN outline and panels
  \draw[very thick] (0,0) rectangle (8,8);
  \node[font=\bfseries,anchor=south] at (4,8.05) {PACMAN 1v5 (top view)};
  \fill[gray!25] (0,0) rectangle (0.35,8);
  \node[font=\scriptsize,rotate=90] at (0.18,3.2) {front panel};
  \fill[gray!25] (7.65,0) rectangle (8,8);
  \node[font=\scriptsize,rotate=90] at (7.83,4) {back panel};
  % front panel connectors, top to bottom
  \node[jack] (sma1) at (0.18,7.3) {};  \node[sig] at (0.5,7.3) {SMA J1001 (P)};
  \node[jack] (sma2) at (0.18,6.6) {};  \node[sig] at (0.5,6.6) {SMA J14 (N)};
  \node[jack] (trig) at (0.18,5.4) {};  \node[sig] at (0.5,5.4) {LEMO TRIG};
  \node[jack] (pps)  at (0.18,4.7) {};  \node[sig] at (0.5,4.7) {LEMO PPS};
  \node[port] (eth)  at (0.18,1.8) {};  \node[sig] at (0.55,1.8) {Ethernet};
  \node[port] (pwr)  at (0.18,0.8) {};  \node[sig] at (0.55,0.8) {power};
  % inside
  \node[box,minimum width=12mm] (j5) at (4.2,7.1) {J5};
  \node[font=\scriptsize,align=center,below=1mm of j5] {ASIC protection jumper\\(removed for checkout)};
  \node[box,minimum width=28mm,minimum height=16mm] (trenz) at (4.2,4.2) {Trenz TE0720\\Linux + firmware};
  \node[box,dashed,minimum width=16mm] (sd) at (4.2,2.0) {microSD};
  \node[tbc,below=1pt of sd] {placement to confirm};
  \draw[thick,fill=white] (7.3,0.3) rectangle (7.65,7.7);
  \node[font=\scriptsize,rotate=90] at (7.05,4) {tile connector (10 tiles)};
  % ---- loopback card
  \draw[very thick,fill=gray!5] (8,0) rectangle (14.4,8);
  \node[font=\bfseries,anchor=south] at (11.2,8.05) {loopback card (Tile 1 shown; $\times$10 tiles)};
  % data loopback
  \draw[ublnk] (8,7.45) -- (12.6,7.45) -- (12.6,6.95) -- (8,6.95);
  \node[sig,anchor=south west] at (8.1,7.45) {TILE1\_POSI\_0..3 (TX)};
  \node[sig,anchor=north west] at (8.1,6.95) {TILE1\_PISO\_0..3 (RX)};
  % timing loopback
  \draw[ublnk,dashed] (8,5.75) -- (12.6,5.75) -- (12.6,5.25) -- (8,5.25);
  \node[sig,anchor=south west] at (8.1,5.75) {TILE1\_CLK, \_RST\_SYNC, \_EXT\_TRIG};
  \node[sig,anchor=north west] at (8.1,5.25) {to an RX input (to be defined)};
  % analog monitor resistor
  \draw[thick] (8,4.15) -- (12.6,4.15) to[R,l=$R_\mathrm{mon}$] (12.6,3.25) -- (8,3.25);
  \node[sig,anchor=south west] at (8.1,4.15) {TILE1\_TESTP};
  \node[sig,anchor=north west] at (8.1,3.25) {TILE1\_TESTN (analog monitor)};
  % loads
  \draw[thick] (8,2.45) -- (12.6,2.45) to[R,l=$R_\mathrm{A}$] (12.6,1.55) node[ground]{};
  \node[sig,anchor=south west] at (8.1,2.45) {TILE1\_VDDA};
  \draw[thick] (8,1.45) -- (10.8,1.45) to[R,l=$R_\mathrm{D}$] (10.8,0.55) node[ground]{};
  \node[sig,anchor=south west] at (8.1,1.45) {TILE1\_VDDD};
  % ---- external instruments
  \node[box,minimum width=26mm,minimum height=12mm] (dmm) at (-4.0,6.95) {DMM\\(across the SMAs)};
  \draw[thick,red!70!black] ([yshift=2mm]dmm.east) -- ++(1.4,0) |- (sma1.west);
  \draw[thick]              ([yshift=-2mm]dmm.east) -- ++(1.9,0) |- (sma2.west);
  \node[box,minimum width=26mm,minimum height=12mm] (fg) at (-4.0,5.05) {function\\generator};
  \coordinate (tee) at (-1.2,5.05);
  \draw[thick] (fg.east) -- (tee);
  \fill (tee) circle (2pt); \node[font=\scriptsize,anchor=east] at ([xshift=-1mm,yshift=2.5mm]tee) {tee};
  \draw[thick] (tee) |- (trig.west);
  \draw[thick] (tee) |- (pps.west);
  \node[box,minimum width=26mm,minimum height=12mm] (pc) at (-4.0,1.8) {operator PC\\(ssh)};
  \draw[thick,{Latex[length=2mm]}-{Latex[length=2mm]}] (pc.east) -- (eth.west);
  \node[box,minimum width=26mm,minimum height=8mm] (ps) at (-4.0,0.8) {power supply};
  \draw[thick] (ps.east) -- (pwr.west);
\end{circuitikz}}
\caption{Checkout setup (physical view; connector placement to be confirmed).  All loopbacks are on the
  loopback card; Tile 1's connections are shown with their signal names, and the card repeats them for all ten
  tiles.  Dashed: routing still to be defined.}
\end{figure}

\begin{center}
\begin{tabularx}{\linewidth}{@{}lL@{}}
\toprule
Item & Role \\
\midrule
Loopback card & On the back-panel tile connector, for all ten tiles: each tile's data outputs (POSI) loop back to
  its data inputs (PISO); its clock, sync and trigger outputs loop back to an RX input; a resistor closes its
  analog monitor pair (TESTP/TESTN); load resistors on VDDA and VDDD let the tile supplies be swept under
  load. \\
DMM & Front-panel SMA readings, entered by the operator when prompted. \\
Function generator & One output, split with a tee: the same square wave on the TRIG and PPS LEMO inputs. \\
Operator PC & \texttt{ssh} session to the PACMAN; keeps the result files. \\
\bottomrule
\end{tabularx}
\end{center}

\begin{itemize}
\item Two versions of the loopback card are needed for different tests, so the checkout is split into one run
  per card, with the PACMAN powered down while the cards are swapped.  Which tests go on which card:
  \tbd{lab day}.
\end{itemize}

\section{The utilities}

\subsection{Common conventions}
\begin{itemize}
\item One line per result: \texttt{<VERDICT> <name>---- <value>}, then a final \texttt{VERDICT} line.
\item Result verdicts: \texttt{PASS}, \texttt{FAIL}, \texttt{SKIP} (not done, with the reason),
  \texttt{INFO} (recorded, not judged).  \texttt{pacman\_hwcheck} also prints \texttt{DATA} lines (sweeps).
\item Run verdict and exit code:
  \begin{center}
  \begin{tabular}{@{}lll@{}}
  \toprule
  Run verdict & Meaning & Exit code \\
  \midrule
  \texttt{PASS} & every check passed, nothing skipped & 0 \\
  \texttt{FAIL} & at least one check failed & 1 \\
  \texttt{INCOMPLETE} & no failure, but something was skipped & 2 \\
  (refused / usage) & bad options, or J5 installed & 3 \\
  \bottomrule
  \end{tabular}
  \end{center}
\item Only a \texttt{PASS} run qualifies a card.
\end{itemize}

\subsection{\texttt{pacman\_lnxcheck}: embedded Linux}

\begin{center}
\begin{tabularx}{\linewidth}{@{}lL@{}}
\toprule
Section & Checks \\
\midrule
System & CPU count, memory size, kernel version. \\
On-chip monitors & Die temperature; six Zynq supplies against nominal. \\
Ethernet & Link speed and duplex (1000/full), MAC address, IP address; ping with small and full-size packets, no
  loss; TFTP transfer of a file twice, compared; interface error counters unchanged. \\
QSPI flash & Each partition read three times; the last two reads must agree.  A difference on the first read
  after boot is reported (INFO), not failed. \\
eMMC & First 64 MB read twice, compared; read rate reported. \\
microSD & Card present, partitions, boot image checksum.  Optional \opt{sd\_write}: a 16 MB test file written,
  read back and removed. \\
\bottomrule
\end{tabularx}
\end{center}
Read-only unless \opt{sd\_write} is given.  Options: \opt{server} (TFTP server; default the default gateway),
\opt{tftp\_file}, \opt{no\_tftp}, \opt{sd\_write}.

\subsection{\texttt{pacman\_hwcheck}: PACMAN hardware}

Run sets execute in this fixed order; \opt{full\_checkout} runs all of them and stops at the first failing
section (\opt{continue} keeps going).

\begin{center}
\begin{tabularx}{\linewidth}{@{}lL>{\raggedright\arraybackslash}p{2.1cm}@{}}
\toprule
Run set & Checks & Operator \\
\midrule
\opt{read\_only} & Firmware and hardware versions, build information, J5 state, board and tile monitors, data
  path status.  Allowed with J5 installed. & -- \\
\opt{blink} & Trenz and PACMAN LEDs blink in turn. & watches, confirms \\
\opt{power} & Board rails against nominal; tile power off, then on.  Each tile in turn: VDDA and VDDD swept over
  the DAC range into the card's loads (floor, gain, swing, monotonic, load), other tiles stay off.  Board-rail
  test points surveyed with the DMM. & DMM readings \\
\opt{loopback} & Each tile in turn, only that tile enabled: random data on its four UARTs through the loopback
  card, words received, bit errors, receiver error counters; no stray data on other UARTs. & -- \\
\opt{timing\_inputs} & Rates on the TRIG and PPS inputs, counted by the firmware over 1 s. & -- \\
\opt{analog\_io} & Front-panel analog MUX stepped through tiles 1--10, test DAC, off; a DMM reading at the SMA
  jacks for each. & DMM readings \\
\bottomrule
\end{tabularx}
\end{center}

\begin{itemize}
\item Operator steps start with a prompt saying what comes next (\texttt{Ready? [y = go, s = skip]}).  At any
  prompt, \texttt{s} skips the step; it is recorded as SKIP, never as a value.
\item Modifiers: \opt{tile\_mask} (tiles present on the fixture), \opt{skip\_interactive} (no operator
  steps), \opt{continue}; development: \opt{internal\_loopback}, \opt{no\_i2c}, \opt{sim\_i2c},
  \opt{ber\_bits}, \opt{analog\_io\_mode}.
\end{itemize}

\needspace{12\baselineskip}
\section{Coverage and acceptance limits}

\begin{longtable}{@{}>{\raggedright\arraybackslash}p{0.25\linewidth}>{\raggedright\arraybackslash}p{0.44\linewidth}>{\raggedright\arraybackslash}p{0.25\linewidth}@{}}
\toprule
Feature & Acceptance & Status \\
\midrule
\endhead
Board rails 3V6, 3V3, 3V0 & monitor within 5\% of nominal; DMM within 5\% of monitor; 3V0 $\le$ 50 mV when off &
  limits set; DMM test points \tbd{} \\
Tile supplies VDDA, VDDD & floor and gain (157 mV per 0x1000 DAC codes, $\pm$5\%), swing $\ge$ 1 V, monotonic,
  current under load, other tiles $\le$ 50 mV & provisional (one board) \\
Tile UARTs (40) & every word received, zero bit errors, no stray words & set \\
TRIG, PPS inputs & rate matches the generator & rates recorded; limits \tbd{} \\
Front-panel analog MUX & expected SMA readings & readings recorded; values \tbd{} \\
LEDs & operator confirms & set \\
Linux, Ethernet, storage & as in the \texttt{pacman\_lnxcheck} table & set \\
\bottomrule
\end{longtable}

\section{Status and untested features}\label{sec:status}

\begin{center}
\begin{tabularx}{\linewidth}{@{}lL@{}}
\toprule
Needed for normal operation, not yet tested & Plan \\
\midrule
Tile CLK, TRIG, SYNC outputs & Loopback to an RX input on a separate connector (fixture \tbd{}). \\
Analog monitor & Known resistor across the input, measured through the front panel. \\
TRIG / PPS acceptance & Function generator at the lab; then rate limits. \\
Full ten-tile run & Loopback card covering all ten tiles. \\
microSD checks & First run on an SD-booted card. \\
\bottomrule
\end{tabularx}
\end{center}

\noindent Not needed for normal operation (not tested): onboard ADC, RTD temperature probe, QSPI flash contents, SFP/CDR
timing path.

\section{Shifter's guide: hardware checkout}

For an operator who runs the checkout the standard way.  Items marked \tbd{\ldots} are still to be filled in.

\subsection*{You need}
\begin{itemize}
\item The PACMAN card with its microSD card (production image), and its serial number \tbd{where printed}.
\item Loopback / load card, DMM, function generator with two outputs, Ethernet cable, PACMAN power supply.
\item A PC on the same network, with \texttt{ssh}.
\end{itemize}

\subsection*{Prepare}
\begin{enumerate}
\item Power off.  Remove J5 (ASIC protection jumper).  Make sure no ASICs are connected.
\item Install the loopback / load card on the tile connectors.
\item Connect the function generator to the TRIG and PPS LEMO inputs: square wave, \tbd{amplitude},
  \tbd{frequency}.
\item Connect Ethernet; power on.  Wait about \tbd{N} s for Linux to boot.
\item Log in: \texttt{ssh \tbd{user}@\tbd{PACMAN address}}.
\item Stop the PACMAN servers:  \texttt{/etc/init.d/pacman\_server stop}
\end{enumerate}

\subsection*{Run}
\begin{enumerate}
\item Linux checkout (a minute or less):\\
  \texttt{pacman\_lnxcheck -{}-server \tbd{TFTP server} | tee lnxcheck-<serial>.txt}
\item Hardware checkout (guided; about \tbd{N} minutes):\\
  \texttt{pacman\_hwcheck -{}-full\_checkout | tee hwcheck-<serial>.txt}
\item Answer the prompts:
  \begin{itemize}
  \item \texttt{Next: \ldots{} Ready?}: get ready for what it describes (watch the LEDs, put the DMM probes on
    the named points), then answer \texttt{y}.
  \item Yes/no questions: answer what you saw.
  \item Readings: type the DMM value in mV.
  \item Can't do a step?  Answer \texttt{s}.  It is recorded as skipped; never type a made-up value.
  \end{itemize}
\end{enumerate}

\subsection*{Read the result}
\begin{itemize}
\item The last line of each run is the verdict.
\item \texttt{PASS} from both: the card passes.
\item \texttt{FAIL}: the lines marked \texttt{FAIL} say what failed.  The hardware checkout stops at the first
  failing section.  Do not repeat blindly: record, and report to \tbd{contact}.
\item \texttt{INCOMPLETE}: something was skipped (lines marked \texttt{SKIP} give the reason).  Fix the setup
  and run again.
\item \texttt{REFUSED} with J5 mentioned: J5 is installed; power off, remove it, start again.
\end{itemize}

\subsection*{Finish}
\begin{enumerate}
\item Copy both result files to \tbd{results location}, named by serial number.
\item Restart the servers (\texttt{/etc/init.d/pacman\_server start}) or reboot.
\item Power off; remove the fixtures.  Install J5 before the card goes to a production system or ASICs.
\item Record the card's status in \tbd{tracking sheet}.
\end{enumerate}

\end{document}
